% ============================================================
% PHẦN 6 (Slide 51--60): FastGS-lite -- kết quả thực nghiệm
% ============================================================

% --- Slide 51 ---
\begin{frame}{FastGS-lite: từ 24GB RTX 4090 xuống Colab T4 miễn phí}
\begin{itemize}
    \item \textbf{fastgs-lite} là bản fork giữ nguyên thuật toán FastGS gốc, chỉ đóng gói lại để chạy trên phần cứng phổ thông (Colab T4 free tier).
    \item Không thay đổi công thức toán học hay pipeline huấn luyện -- chỉ thay đổi cách triển khai/vận hành.
\end{itemize}
\vspace{0.3em}
\footnotesize
\begin{table}
\centering
\begin{tabular}{lll}
\toprule
\textbf{Tiêu chí} & \textbf{FastGS gốc} & \textbf{fastgs-lite} \\
\midrule
Phần cứng mục tiêu & 24GB RTX 4090 & T4 ($\sim$15GB VRAM) + 12.7GB RAM \\
CUDA extension & git submodule & vendor sẵn trong repo (\texttt{git clone}) \\
Entry point & train.py + shell script & train.py + notebook có instrumentation \\
Tín hiệu tiến trình & L1 + EMA & Score tổng hợp, log mỗi 1000 iter \\
\bottomrule
\end{tabular}
\end{table}
\end{frame}

% --- Slide 52 ---
\begin{frame}{Ràng buộc phần cứng thực tế}
\begin{itemize}
    \item Colab T4 free tier: \textbf{14.56 GB VRAM}, nhưng RAM hệ thống (host) chỉ \textbf{12.7 GB}.
    \item Trực giác thông thường: nghẽn cổ chai sẽ ở VRAM (GPU). \textbf{Thực tế ngược lại}.
    \item VRAM peak đo được chỉ $\sim$1.16GB -- khoảng \textbf{8\%} dung lượng khả dụng.
    \item RAM peak đo được $\sim$9.59GB -- khoảng \textbf{91\%} giới hạn mềm 10.5GB của pipeline.
    \item[$\Rightarrow$] Phát hiện quan trọng, sẽ nhắc lại ở phần \emph{Quan sát} cuối bài: \textbf{RAM hệ thống, không phải VRAM, là ràng buộc thật sự}.
\end{itemize}
\vspace{0.5em}
\centering
\footnotesize
\begin{tabular}{lcc}
\toprule
 & VRAM (14.56GB) & RAM (12.7GB, soft limit 10.5GB) \\
\midrule
Peak sử dụng & 1.16GB (8\%) & 9.59GB (91\% soft limit) \\
\bottomrule
\end{tabular}
\end{frame}

% --- Slide 53 ---
\begin{frame}{Kiến trúc notebook: 7 cell $\leftrightarrow$ 7 module}
\footnotesize
\begin{table}
\centering
\begin{tabular}{cll}
\toprule
\textbf{Cell} & \textbf{Module} & \textbf{Chức năng} \\
\midrule
1 & \texttt{env.py} & Cài dependency, kiểm tra GPU \\
2 & \texttt{config.py} & Một object \texttt{Config} duy nhất cho toàn pipeline \\
3 & \texttt{data.py} & Tải dữ liệu, liệt kê scene \\
4 & \texttt{run.py} & Smoke test vài trăm iteration \\
5 & \texttt{trainer.py} + \texttt{score.py} & Train toàn bộ scene, log Score + \% tiến trình \\
6 & \texttt{report.py} & Bảng/biểu đồ phân tích kết quả \\
7 & \texttt{submission.py} + \texttt{deliver.py} & Đóng gói \texttt{submission.zip}, kiểm tra format, tự tải về \\
\bottomrule
\end{tabular}
\end{table}
\normalsize
\begin{itemize}
    \item Ánh xạ 1--1 giữa cell notebook và module trong \texttt{pipeline/} giúp dễ debug và tái sử dụng ngoài Colab.
\end{itemize}
\end{frame}

% --- Slide 54 ---
\begin{frame}{Thiết lập thực nghiệm}
\begin{itemize}
    \item 4 scene từ 2 dataset: \textbf{Deep Blending} (\texttt{drjohnson}, \texttt{playroom}) và \textbf{Tanks\&Temples} (\texttt{train}, \texttt{truck}).
    \item Split \texttt{--eval} với \texttt{llffhold=8} (cứ 8 ảnh lấy 1 làm test).
    \item Cấu hình tối ưu giống nhau mọi scene: \texttt{iterations=7000}, \texttt{sh\_degree=3}, \texttt{resolution=2} khi train.
    \item \textbf{Chú ý}: eval lại dùng độ phân giải gốc (\texttt{submission\_resolution=1}) -- sự lệch train/eval này do \emph{protocol cuộc thi}, không phải sơ suất.
\end{itemize}
\vspace{0.3em}
\footnotesize
\centering
\begin{tabular}{lcccc}
\toprule
\textbf{Scene} & \textbf{Cameras} & \textbf{Train views} & \textbf{Test views} & \textbf{Res. eval} \\
\midrule
drjohnson & -- & 7/8 & 1/8 & gốc (res=1) \\
playroom  & -- & 7/8 & 1/8 & gốc (res=1) \\
train     & -- & 7/8 & 1/8 & gốc (res=1) \\
truck     & -- & 7/8 & 1/8 & gốc (res=1) \\
\bottomrule
\end{tabular}
\end{frame}

% --- Slide 55 ---
\begin{frame}{Bảng 1: Chất lượng cuối cùng theo scene}
\footnotesize
\begin{table}
\centering
\begin{tabular}{lccccc}
\toprule
\textbf{Scene} & \textbf{Score} & \textbf{PSNR (dB)} & \textbf{SSIM} & \textbf{LPIPS} & \textbf{\#Gaussian} \\
\midrule
playroom  & 0.8198 & -- & -- & -- & $\sim$129k \\
drjohnson & 0.8027 & 27.41 & -- & -- & $\sim$174k \\
truck     & 0.7482 & -- & -- & -- & $\sim$187k \\
train     & 0.6867 & 19.75 & -- & -- & $\sim$201k \\
\midrule
\textbf{Mean} & \textbf{0.7643} & \textbf{24.46} & & & \\
\bottomrule
\end{tabular}
\end{table}
\begin{itemize}
    \item 2 scene indoor (playroom, drjohnson) vượt 2 scene outdoor (truck, train) dù dùng \textbf{ít Gaussian hơn} (129k--174k so với 187k--201k).
\end{itemize}
\centering
\includegraphics[width=0.55\linewidth]{leaderboard.png}
\captionof{figure}{Leaderboard Score theo scene}
\end{frame}

% --- Slide 56 ---
\begin{frame}{Bảng 2: Động lực huấn luyện}
\footnotesize
\begin{itemize}
    \item Score đo mỗi 1000 iteration trên cả 4 scene.
    \item Hai điểm sụt rõ rệt tại iteration \textbf{3000} và \textbf{6000}: Score rơi xuống $\sim$0.22--0.31.
    \item Trùng khớp chính xác với \texttt{opacity\_reset\_interval=3000}.
    \item[$\Rightarrow$] Đây là \textbf{hiện tượng đo lường} (measurement artifact) do lấy mẫu Score ngay lúc reset opacity, \textbf{không phải} mô hình phân kỳ. Phục hồi hoàn toàn trong dưới 1000 iteration sau đó.
\end{itemize}
\centering
\includegraphics[width=0.48\linewidth]{training.png}
\captionof{figure}{\footnotesize Score theo iteration -- các điểm sụt tại 3000, 6000}
\end{frame}

% --- Slide 57 ---
\begin{frame}{Kết quả định tính: drjohnson (tốt nhất)}
\footnotesize
\begin{itemize}
    \item Score \textbf{0.8027}, PSNR \textbf{27.41dB}. Hình học, tư thế, màu sắc gần như khớp hoàn toàn với ground truth ở view đầu.
    \item Lỗi còn lại: chi tiết tần số cao / phụ thuộc góc nhìn -- chi tiết mảnh (khe tản nhiệt, gáy sách) bị mờ, highlight phản chiếu trên gỗ mahogany gần như biến mất.
    \item Nguyên nhân: các lobe SH bậc cao (view-dependent) có learning rate bị chia thêm \textbf{/20} trong code.
\end{itemize}
\centering
\includegraphics[width=0.5\linewidth]{samples_drjohnson.png}
\captionof{figure}{\footnotesize So sánh render vs. ground truth -- drjohnson}
\end{frame}

% --- Slide 58 ---
\begin{frame}{Kết quả định tính: train (yếu nhất -- vấn đề bầu trời)}
\footnotesize
\begin{itemize}
    \item Score \textbf{0.6867}, PSNR \textbf{19.75dB} -- thấp nhất trong 4 scene.
    \item Thất bại rõ rệt nhất: \textbf{vùng trời}. Ground truth là màu xanh đồng nhất, render lại đầy vệt xám-trắng loang lổ. Đầu máy xe lửa (có texture/parallax) tái tạo tốt.
    \item Vùng trời không texture, ở khoảng cách vô hạn, \textbf{không cho tín hiệu multi-view} nào để loại bỏ Gaussian dư thừa.
    \item[$\Rightarrow$] Giải thích vì sao lỗi rơi vào \textbf{PSNR} chứ không phải \textbf{LPIPS} (perceptual, ít nhạy vùng đồng nhất).
\end{itemize}
\centering
\includegraphics[width=0.5\linewidth]{samples_train.png}
\captionof{figure}{\footnotesize So sánh render vs. ground truth -- train (lỗi vùng trời)}
\end{frame}

% --- Slide 59 ---
\begin{frame}{Bảng 3: Chi phí \& lưu trữ}
\begin{itemize}
    \item Tổng thời gian train 4 scene: \textbf{327.6 giây} ($\sim$5.5 phút), tốc độ trung bình \textbf{85.5 iter/s}.
    \item VRAM đỉnh trung bình chỉ \textbf{0.88GB} (tối đa 1.16GB ở drjohnson) -- rất thấp so với 14.56GB khả dụng.
    \item Kích thước \texttt{point\_cloud.ply} trung bình \textbf{171.6MB} / 4 scene.
    \item Đúng \textbf{248 byte/Gaussian} = kích thước record 3DGS không nén (62 trường float32: 3 vị trí + 3 normal + 3 f\_dc + 45 f\_rest + 1 opacity + 3 scale + 4 rotation).
\end{itemize}
\vspace{0.3em}
\footnotesize
\centering
\begin{tabular}{lcc}
\toprule
\textbf{Chỉ số} & \textbf{Trung bình} & \textbf{Ghi chú} \\
\midrule
Thời gian train (4 scene) & 327.6 s & $\sim$5.5 phút tổng \\
Tốc độ & 85.5 iter/s & \\
VRAM peak & 0.88 GB & max 1.16GB (drjohnson) \\
Kích thước .ply & 171.6 MB & 248 byte/Gaussian \\
\bottomrule
\end{tabular}
\end{frame}

% --- Slide 60 ---
\begin{frame}{Khả năng tái lập \& giới hạn phép đo}
\footnotesize
\texttt{Config(resolution=2, iterations=7000, score\_every=1000,\\
\hspace{2em}eval\_views=6, psnr\_max=30.0, submission\_resolution=1)}
\normalsize
\vspace{0.5em}
\begin{itemize}
    \item \textbf{Giới hạn 1}: đây là \textbf{một lần chạy duy nhất}, không cố định seed, không lặp lại nhiều lần $\Rightarrow$ không có ước lượng phương sai/error bar.
    \item \textbf{Giới hạn 2}: chỉ 4 scene từ 2 dataset -- mẫu hẹp.
    \item \textbf{Giới hạn 3}: ngân sách 7k iteration chỉ bằng \textbf{1/3} lịch trình optimizer thiết kế cho 30k.
    \item \textbf{Giới hạn 4}: bỏ qua hoàn toàn bước prune cuối (\texttt{final\_prune\_fastgs}).
    \item[$\Rightarrow$] \textbf{Kết luận}: Bảng 1 nên xem là \emph{baseline có thể tái lập của repo này}, \textbf{không phải} kết quả benchmark chính thức để so sánh với FastGS/3DGS gốc.
\end{itemize}
\end{frame}
